\section{Introduction}\label{sec:intro}

In the realm of software systems, such as device drivers, file
systems, and network stacks, precise control over the
\emph{data layout} of objects is crucial for compatibility and
performance. Specifically, controlling the composition of objects in memory
on a bit- and byte-level can avoid the need for translation or \emph{deserialisation} 
at the boundaries between on-medium and
in-memory data which frequently arise from interacting with
standardised protocols or hardware.
%
These systems are often implemented in the C language, in part because it offers low-level features to give fine-grained control over data layout. 
Unfortunately, to maintain good performance, the C programmer must throw away the conceptual abstraction of the data type, and instead focus on the low-level details of bits and bytes. This low-level code contains many subtle bit-twiddling
operations which, apart from being difficult to manually verify, are
also tedious and error-prone to implement.

% 
In Linux, a \cc{dev_t} type is used to represent device numbers. It is \cc{typedef}ed to
a 32-bit unsigned integer (\cc{u32}). It is composed of two parts: a major number and a minor
number. Within the 32 bits, 12 bits are set aside for the major number, and the remaining
20 bits for the minor number. In C, to construct a \cc{dev_t} or to fetch the device
numbers, users need to use a some macros, which are implemented with convoluted bit twiddling:
\begin{lstlisting}[language=C]
#define MINORBITS    20
#define MINORMASK    ((1U << MINORBITS) - 1)
#define MAJOR(dev)   ((unsigned int) ((dev) >> MINORBITS)) // get major
#define MINOR(dev)   ((unsigned int) ((dev) & MINORMASK))  // get minor
#define MKDEV(ma,mi) (((ma) << MINORBITS) | (mi))          // construct dev_t
\end{lstlisting}
Logically, the \cc{dev_t} type is simply a product of two integers.
To implement it in a high-level language, users would typically define it as a record
type with two integer fields. For example, in \cogent, it may be defined as:\footnote{
As an opening example, we show pseudo-\cogent code here to omit the irrelevant \cogent
intricacies.}

\begin{lstlisting}[language=Cogent]
type Dev_t = { major : Int, minor : Int }
\end{lstlisting}
and to access the fields of such an object \cg{dev},
the normal member operation (\cg{dev.major} and \cg{dev.minor}) can be used.
With this abtract view of the type, which facilitates simpler operations on this type,
the user unfortunately has to hand over control of the type's layout to the
\cogent compiler. The user then needs to study the compiler's code generation convention
and write data format conversion functions to related the compiled \cg{Dev_t} type with
the native C type \cc{dev_t}. In fact, the user does not have to make this compromise. With
\dargent, the user can have both the abstraction and control over types' memory
layouts. By specifying a \dargent layout description, the \cogent compiler can
generate C code with the specified layout:
\begin{lstlisting}[language=Cogent]
type Dev_t = { major : Int, minor : Int }
  layout record { major : 12b, minor : 20b }
\end{lstlisting}
This is transparent to the functional smenatics
of the \cogent program. Namely, the \cg{dev.major} operation continues to work regardless
of the underlying layout of \cg{dev}.





To maintain the higher-level structure of a program without sacrificing
performance, we want to use a language with a high-level semantics, but with facilities for
specifying the low-level memory layout of heap-allocated objects. While most
high-level languages use fixed heap layouts, there has been recent
progress on language-based support for user-defined memory
layouts~\cite{Cronburg:2019,Vollmer:2019}. These languages, however, do
not provide verified correctness of their translations.

In this paper, we present \Dargent, a language for describing data
layouts of high-level algebraic datatypes along with a data
refinement framework for \emph{automatically verifying} the correctness of the compiled C code
with respect to those layout descriptions. We
build on the \cogent language and refinement framework~\citep{Cogent_JFP_2021}.
\cogent is designed for the implementation of high-assurance low-level systems components as pure
mathematical functions operating on algebraic data types.
\cogent's certifying compiler co-generates
a C program and Isabelle/HOL~\citep{Nipkow:2002} theorems witnessing a proof
that the C program refines an Isabelle/HOL embedding of the \cogent source program~\cite{Rizkallah:itp16}. 

While there is a very long line of prior work on \emph{data description
languages}~\citep{FisherG05, McCannC00, Back02, WangG11, everparse,
Ye_Delaware_19, vanGeest_Swierstra_17, Slind_21}, and the data layout
descriptions used in \dargent do indeed look similar to those used in such
languages, there is a fundamental difference: These languages are
designed for synthesising data (de)serialisation functions (also sometimes
referred to as data marshalling/unmarshalling functions, encoders/decoders, 
or parsers and pretty-printers for low-level data),
which convert data stored in a low-level, sequential format in some storage
medium to a high-level, structured representation in memory and vice versa. 
They are primarily used for the interaction and communication between
different programming languages (e.g.\ a foreign function interface)
or systems (e.g.\ data transmission over the network). In this context, code to transform back-and-forth
between the two representations is still necessary.

\dargent, on the other hand, is intended to solve a different
problem. The \dargent data layout descriptions grant programmers
the ability to dictate to the \cogent compiler how it should lay out the algebraic data
types used by the \cogent program itself. The compiler generates code that works directly 
with data laid out according to the programmer's specifications, as well as Isabelle/HOL  
proofs showing that it has done so correctly. 

Depending on the application, \dargent therefore can either eliminate entirely
or reduce the need for data (de)serialisation code, be it manually
written or automatically derived, when interacting with the external world.
Because algebraic data types can be represented directly in their binary data formats with \dargent,  the 
programmer does not need
to decode raw data first into some other in-memory representation in order to operate on it as a data type. Eliminating these (de)serialisation steps naturally results in more concise and readable code, better performance, and easier informal and 
formal reasoning.

This additional power in expressiveness can also be
used to improve the performance of the compiled \cogent code, 
e.g.\ by having smaller memory footprints or a specialised memory
layout optimal for the underlying architecture,
independent of the interoperation between languages or systems.
It also enables \cogent programmers to directly write code that is binary-compatible
with native C programs.

\cogent is readily amenable to an extension for prescribing 
data layouts and fine-tuning the compilation of algebraic data types 
by virtue of its lack of a runtime system
and direct compilation to C. Before \dargent, it simply adopted
the layout conventions of the
underlying C compiler. The introduction of \dargent into the
framework enabled improvements to
some outstanding inefficiencies in the prior design, such as
a reduced reliance on deserialisation code within file system
implementations~\citep{Amani:2016} and directly representing device
register formats as data types.
%%
Furthermore, we have extended the \cogent compiler so as to preserve the
benefits of \cogent's high-level type system and semantics. The upshot
is our compiler automatically translates read and write operations on
heap-allocated objects to take account of their particular data
layout. The translation's correctness is guaranteed by the enhanced
data refinement theorem. 
This paper describes, to the best of our knowledge, the first framework that
is able to leverage data layout specifications for generating bit-level accessors
with formal correctness guarantees. %Thanks to \cogent's uniqueness type system, our framework also supports mutable data. 
%This is the first \emph{certifying} framework to be extended with data layout descriptions supporting mutable data, thanks to \cogent's
%uniqueness type system which allows for efficient destructive
%updates.

\subsection*{Contributions}
This paper realises the vision set out previously~\citep{Isola:2018}.
We make the following contributions:
\begin{itemize}
  \item The design and implementation of \Dargent
    (\autoref{sec:dargent}), a \emph{data layout} language for
    controlling the memory layout of algebraic data types, down to the
    bit level. We formalise a core calculus and its static semantics
    in Isabelle/HOL~\citep{Nipkow:2002}, and discuss the compilation process;
  \item An extension to the \cogent verification framework 
    (\autoref{sec:verification}) to \emph{automatically verify} the
    translation of high-level read/write accesses (known as
    \emph{getters} and \emph{setters}) to explicit offsets within a
    well-defined memory region;
  \item An extended suite of examples (\autoref{sec:applications}
    and \autoref{sec:timerdriver}) demonstrating the utility of \Dargent in
    the context of device drivers.
\end{itemize}
All the results described in this paper, including the case studies we
present, are associated with formal proofs in Isabelle/HOL.
These materials are packaged as a virtual machine image, which is publicly
available~\citep{dargent-artefact}. It is derived from a development branch of
the \cogent project repository (\url{https://github.com/au-ts/cogent}).

